% Proposed replacement for the paragraph immediately following
% equations bloch:eq:goldenP3 and bloch:eq:sgP3 in
% Analysis/Polylogarithms/docs/manuscript/chapters/03-algebraic.tex.
% At the audited revision its last sentence (lines 428--429) asserts
% that the weight-6 golden ladder becomes a nonzero zeta(6) multiple
% under P_m.  This is incompatible with the chapter's definition of
% P_m: P_{2j}(x) is the imaginary part of a real expression for 0<x<1.

Both displayed weight-three identities are exact.
Compare~\eqref{bloch:eq:goldenP3} with the ordinary formula
$15\Li_3(\phi^{-2})=10\ln^3\phi-2\pi^2\ln\phi+12\zeta(3)$:
in these examples, passing to $P_3$ cancels the lower-weight
logarithmic terms and leaves the displayed $\zeta(3)$ multiple.
At higher odd weights the corresponding lower-weight terms in the
definition of $P_m$ must likewise be included, with the proof status
of each resulting identity stated separately.  At even weight this
particular modified polylogarithm satisfies
$P_{2j}(x)=0$ for every $0<x<1$.  Therefore the nonzero
$\zeta(6)$ and $\zeta(8)$ multiples in the ordinary golden formulas
are not evaluations of $P_6$ and $P_8$ at those real arguments.
Those even-weight identities remain ordinary-polylogarithm formulas,
unless a different real normalization is explicitly defined.
